% ---- pg-166 （书页 154） ----

\noindent (a) $X_1$ 和 $X_2$

\noindent (b) $X_2$ 和 $X_3$

\noindent (c) $(X_1, X_2)$ 和 $X_3$

\noindent (d) $\dfrac{X_1 + X_2}{2}$ 和 $X_3$

\noindent (e) $X_2$ 和 $X_2 - \dfrac{5}{2}X_1 - X_3$

\noindent 4.4\quad 设 $\mathbf{X}$ 服从 $\boldsymbol{\mu}' = [2, -3, 1]$ 和
\[
\boldsymbol{\Sigma} = \begin{bmatrix} 1 & 1 & 1 \\ 1 & 3 & 2 \\ 1 & 2 & 2 \end{bmatrix}
\]
的 $N_3(\boldsymbol{\mu}, \boldsymbol{\Sigma})$,

\noindent (a) 求 $3X_1 - 2X_2 + X_3$ 的分布.

\noindent (b) 如果必要, 重新给变量标号, 并求一个 $2 \times 1$ 向量 $\mathbf{a}$, 使得 $X_2$ 与 $X_2 - \mathbf{a}'\begin{bmatrix} X_1 \\ X_3 \end{bmatrix}$ 相互独立.

\noindent 4.5\quad 对以下每一个概念作出详细说明.

\noindent (a) 对练习 4.2 中的联合分布给定 $X_2 = x_2$ 时 $X_1$ 的条件分布.

\noindent (b) 对练习 4.3 中的联合分布给定 $X_1 = x_1$ 和 $X_3 = x_3$ 时 $X_2$ 的条件分布.

\noindent (c) 对练习 4.4 中的联合分布给定 $X_1 = x_1$ 和 $X_2 = x_2$ 时 $X_3$ 的条件分布.

\noindent 4.6\quad 设 $\mathbf{X}$ 服从 $N_3(\boldsymbol{\mu}, \boldsymbol{\Sigma})$ 分布, 其中 $\boldsymbol{\mu}' = [1, -1, 2]$,
\[
\boldsymbol{\Sigma} = \begin{bmatrix} 4 & 0 & -1 \\ 0 & 5 & 0 \\ -1 & 0 & 2 \end{bmatrix}
\]
下列各对随机变量中哪几对相互独立? 试说明.

\noindent (a) $X_1$ 和 $X_2$

\noindent (b) $X_1$ 和 $X_3$

\noindent (c) $X_2$ 和 $X_3$

\noindent (d) $(X_1, X_3)$ 和 $X_2$

\noindent (e) $X_1$ 和 $X_1 + 3X_2 - 2X_3$

\noindent 4.7\quad 参看练习 4.6 详细说明以下每一个概念.

\noindent (a) 给定 $X_3 = x_3$ 时 $X_1$ 的条件分布.

\noindent (b) 给定 $X_2 = x_2$ 和 $X_3 = x_3$ 时 $X_1$ 的条件分布.

\noindent 4.8\quad (一个有正态边缘的非正态二元分布的例子)$X_1$ 服从 $N(0, 1)$, 且
\[
X_2 = \begin{cases} -X_1 & \text{若 } -1 \leq X_1 \leq 1 \\ X_1 & \text{其他} \end{cases}
\]
证明下列命题:

\noindent (a) $X_2$ 也服从 $N(0, 1)$ 分布.

\noindent (b) $X_1$ 和 $X_2$ 不服从二元正态分布.

\noindent 提示:

\noindent (a) 由于 $X_1$ 服从 $N(0, 1)$, 对任意 $x$ 有 $P[-1 < X_1 \leq x] = P[-x \leq X_1 < 1]$. 当 $-1 < x_2 < 1$
